% Edge cases that the two example manuscripts do not cover.
\documentclass[sn-mathphys-ay]{sn-jnl}
\usepackage{amsmath}
\usepackage{amsthm}
\usepackage{subfig}
\newcommand{\vect}[1]{\mathbf{#1}}
\newcommand{\sample}{Sample text}
\newcommand{\figref}[1]{Fig.~\ref{#1}}
\DeclareMathOperator{\diag}{diag}
\newtheorem{lemma}{Lemma}
\newtheorem{corollary}[lemma]{Corollary}

\begin{document}
\title{Edge cases for the \LaTeX\ to Word converter}
\author*[1]{\fnm{Ada} \sur{Lovelace}}\email{ada@example.org}
\author[1,2]{\fnm{Charles} \sur{Babbage}}\email{charles@example.org}
\equalcont{Equal contribution.}
\affil*[1]{\orgname{Analytical Engine Society}, \orgaddress{\city{London}, \country{UK}}}
\affil[2]{\orgname{Difference Engine Lab}, \orgaddress{\city{Cambridge}, \country{UK}}}
\abstract{An abstract that cites \citep{knuth84} and mentions \sample.}
\keywords{edge cases; macros}
\maketitle

\section{Macros and theorems}\label{sec:macros}
\sample{} with a vector $\vect{v}$ and an operator $\diag(x)$.% a trailing comment
Joined line.

\begin{lemma}\label{lem:first}
A lemma with an equation
\begin{equation}
  a = b \tag{$\star$}\label{eq:star}
\end{equation}
\end{lemma}
\begin{corollary}\label{cor:second}
A corollary sharing the lemma counter.
\end{corollary}
See Lemma~\ref{lem:first}, Corollary~\ref{cor:second}, equation~\eqref{eq:star},
\figref{fig:sub} and panel~\ref{fig:sub:b}. An undefined reference \ref{nope}.
A footnote\footnote{Footnote pointing to Section~\ref{sec:macros}.} here.

\begin{gather*}
  x = 1 \\
  y = 2
\end{gather*}
\begin{align}
  p &= q \nonumber \\
  r &= s \label{eq:rs}
\end{align}

\begin{figure}
  \centering
  \subfloat[Left panel\label{fig:sub:a}]{\includegraphics[width=0.4\textwidth]{missing-image}}
  \subfloat[Right panel\label{fig:sub:b}]{\includegraphics[width=0.4\textwidth]{missing-image}}
  \caption{Two sub-floats.}\label{fig:sub}
\end{figure}

\section{Algorithms}
\begin{algorithm}
\caption{Loop}\label{alg:loop}
\begin{algorithmic}[1]
\Procedure{Run}{$n$}
  \For{$i = 1$ to $n$}\label{line:for}
    \State \Call{Work}{$i$} \Comment{do work}
  \EndFor
  \State \Return $n$
\EndProcedure
\end{algorithmic}
\end{algorithm}
Line~\ref{line:for} of Algorithm~\ref{alg:loop}. Cite \citet{lamport94} and \citep{knuth84,lamport94}.

\begin{thebibliography}{2}
\bibitem[{Knuth(1984)}]{knuth84} D. E. Knuth. \newblock \emph{The {\TeX}book}. \newblock Addison-Wesley, 1984.
\bibitem[{Lamport(1994)}]{lamport94} L. Lamport. \newblock \emph{{\LaTeX}: A Document Preparation System}. \newblock Addison-Wesley, 1994.
\end{thebibliography}
\end{document}
